\scanpage{81}
de $\handwritten{\underline}{Q}_p$ \handwritten{)} qui sont donc isomorphes \inserted{\handwritten{``géométriquement'' i.e. dans $\operatorname{Rep}(\underline{H})$}} (mais sans doute non canoniquement) aux $E_{r,1}$ pour $r\geq1$: ils donnent des modules libres de rang 1 canoniques sous les $K_{r,1}$ ...\handwritten{]}

9. Pour $k$ quelconque, on trouve un \deleted{ho} $\handwritten{\otimes}$-homomorphisme canonique \deleted{(à l'aide} défini à isomorphisme unique près (on utilise un choix d'une clôture algébrique $\handwritten{\overline}{k}$ de $k$, mais ce choix est inessentiel ..)
\[
\operatorname{Fcriso}(k)\handwritten{\longrightarrow}\operatorname{Rep}(\handwritten{\underline}{H})
\]
ou encore un homomorphisme de gerbes
\[
\handwritten{\underline}{H}\handwritten{\longrightarrow}\handwritten{\underline}{G}(k)
\]
(rappelons que cet homomorphisme se factorise par le noyau,\deleted{de} soit \deleted{\unsure{G|Gr}} $\handwritten{\underline}{H}(k)$, de $\handwritten{\underline}{G}(k)\handwritten{\to}\handwritten{\underline{\Pi}}_p(k)$, et que pour $k$ parfait on devine que l'homomorphisme obtenu $\handwritten{\underline}{H}\handwritten{\to}\handwritten{\underline}{H}(k)$ est une équivalence ...). Si $\handwritten{\lambda\in}\handwritten{\underline}{Q}$, un F\deleted{(}\handwritten{-}isocristal sur $k$ est dit \deleted{\unsure{homogèn|homogène}} \handwritten{\uline}{de pente} $\handwritten{\lambda}$ si son image dans $\operatorname{Rep}(\handwritten{\underline}{H})$ est de pente $\handwritten{\lambda}$. On voit alors aussitôt qu'un F-isocristal est de pente zéro sss il est ``trivial'' i.e. satisfait aux conditions de 5.2 (ce qui justifie la terminologie qui est introduite, un peu prématurément, à cet endroit), il est effectif sss son image dans \inserted{$\operatorname{Rep}(H)$ l'est,i.e. est à pente $\handwritten{\geq}$ 0.}

Considérons l'homomorphisme de liens associé
\[
H\handwritten{\longrightarrow}G(k)\quad ,
\]
on trouve \unsure{alons|alors} que si $k$ est parfait, l'homomorphisme précédent est \handwritten{\uline}{central}, en d'autres termes $H$ opère sur le foncteur identique de $\operatorname{Fcriso}(k)\handwritten{\to}\operatorname{Rep}(\handwritten{\underline}{G}(k))$, de sorte que tout F-isocristal \inserted{$(E,F_E)$} se décompose en ses \handwritten{``}\handwritten{\uline}{composantes isopentiques}\handwritten{''} $E_{\handwritten{\lambda}}$, $\handwritten{\lambda\in}\handwritten{\underline}{Q}$. $E$ est ``trivial'' resp. effectif sss $E_{\handwritten{\lambda}}=0$ pour $\handwritten{\lambda\ne}$ 0 (resp. $\handwritten{\lambda<0}$).

Notons que le \inserted{F-iso}cristal de Tate $K(\unsure{1|l})$ est de pente\handwritten{-}1. Donc (pour $k$ parfait) tout F\deleted{(}\handwritten{-}isocristal se met de façon unique sous la forme
\[
E=\handwritten{\sum_{i\in\mathbb{Z}}}E_i(-i)\quad ,
\]
où $E_i$ est à composantes isopentiques de pente $0\handwritten{<}\,\inserted{\handwritten{i}}\handwritten{\leq}1$. D'ailleurs,
% archive sheet number: 80
